\section{Preprocessing}

The preprocessing pipeline is structured around the four classical
data-mining stages from the CS313 curriculum: \textbf{cleaning},
\textbf{integration}, \textbf{transformation}, and \textbf{reduction}.
The full implementation lives in \texttt{preprocessing/preprocess.py};
the input is the consolidated 131M-row interaction parquet, the output
is \texttt{kt4\_preprocessed.parquet} (131.0M rows, 30 columns, 2.4 GB).

\subsection{Cleaning}

\textbf{Missing values.} As established during EDA, all four columns
with missing values are structurally missing — present 100\% of the
time when applicable. No imputation is applied.

\textbf{Duplicates.} 461{,}237 exact duplicate rows (0.35\%) were
detected and removed. These almost certainly originate from
client-side retry logic (network blips re-sending the same event).
Removal is lossless.

\textbf{Consistency.} Programmatic checks confirmed that
\texttt{user\_answer} values are always in \{a, b, c, d\},
\texttt{platform} is always in \{mobile, web\} (or null where
structurally expected), no timestamps are non-positive, and per-user
event timestamps are monotonically non-decreasing in a 10K-user
sample. \textbf{Zero violations} found.

\textbf{Outliers (cursor\_time).} Media-action playback positions
range up to 11.6M ms (~3.2 hours) — clearly accidental
left-running sessions. Using the IQR rule on the non-null subset
(Q1 = 0, Q3 = 17{,}650 ms, upper bound = 44{,}125 ms), \textbf{2.68M
extreme \texttt{cursor\_time} values were set to NaN}. The rows
themselves are preserved — only the unreliable measurement is
cleared, so the subsequent min-max normalization sees a meaningful
0--44s range instead of being crushed by 3-hour outliers.

\subsection{Integration}

We derived \texttt{item\_type} from the prefix of \texttt{item\_id}
(\texttt{q}/\texttt{b}/\texttt{l}/\texttt{e}/\texttt{p}/\texttt{c})
and used it to drive joins with the metadata files:

\begin{itemize}
\setlength\itemsep{1pt}
\item \textbf{Questions join}: 28{,}142{,}959 question-related rows
  matched against \texttt{questions.csv} at \textbf{100\%} success;
  added \texttt{bundle\_id}, \texttt{correct\_answer}, \texttt{part},
  \texttt{tags}.
\item \textbf{Lectures join}: 5{,}009{,}098 lecture-related rows
  matched against \texttt{lectures.csv} at \textbf{100\%} success;
  added \texttt{lecture\_part}, \texttt{lecture\_tags},
  \texttt{lecture\_video\_length}.
\item \textbf{Correctness}: \texttt{is\_correct =
  (user\_answer == correct\_answer)} computed for every
  \texttt{respond} action — 23{,}308{,}702 labeled rows in total,
  with \textbf{56.87\% correct}. This single column is the foundation
  for every downstream prediction task.
\item \textbf{Redundancy}: \texttt{explanation\_id} was found to
  always equal \texttt{bundle\_id} for every question and was therefore
  not added; \texttt{bundle\_id} and \texttt{item\_type} were kept
  even though derivable, because they avoid repeated joins downstream.
\end{itemize}

\subsection{Transformation}

\textbf{Attribute construction}: from \texttt{timestamp} we derived
\texttt{hour}, \texttt{day\_of\_week}, \texttt{date},
\texttt{time\_since\_prev} (per-user inter-event gap in ms), and
\texttt{action\_seq} (per-user 0-indexed action number).

\textbf{Normalization}: \texttt{cursor\_time\_normalized} is
\texttt{cursor\_time} min-max-scaled to $[0, 1]$ over the cleaned
0--44s range; \texttt{log\_time\_since\_prev = log(1 +
time\_since\_prev)} tames the extreme right skew (median 3 s,
max 446 days) into a distribution amenable to downstream models.

\textbf{Discretization}: \texttt{time\_of\_day} bins \texttt{hour}
into \{night, morning, afternoon, evening\}; \texttt{is\_weekend}
binary-encodes Sat/Sun.

\textbf{Encoding}: integer label encoding for the four categorical
columns (\texttt{action\_type}, \texttt{source}, \texttt{platform},
\texttt{item\_type}); the original strings are preserved alongside.

\subsection{Reduction (deferred)}

The full preprocessed dataset is saved without reduction so the team
retains flexibility for downstream tasks. The report documents six
reduction strategies for downstream use — user sampling, stratified
sampling, attribute subset selection, aggregation, cold-start
filtering, and PCA — to be invoked task-by-task rather than baked
into the canonical preprocessed file.

\begin{figure}[H]
\centering
\includegraphics[width=0.78\linewidth]{preprocessing_summary.png}
\caption{Preprocessing pipeline summary: input (131.4M, 8 cols) $\to$ output (131.0M, 30 cols).}
\end{figure}
